% !TeX program = xelatex
% =============================================================
%  Reporte final — Práctica 2: Emulador básico de un procesador
%  Copia modificada de docs/practica/final/plantillas/Reporte_LaTeX/main.tex
%  (mismo estilo UV: colores, tipografía, encabezados, pies).
%
%  Compilar:  xelatex reporte_final.tex   (2 veces, para el índice)
%  En Linux sin TeX instalado: subir la carpeta a Overleaf y compilar
%  con XeLaTeX. Las imágenes viven en imagenes/.
%
%  NOTA: la portada y la anteportada de la plantilla UV se agregan
%  después desde el editor (las tiene el equipo en PDF).
% =============================================================
\documentclass[11pt,letterpaper,oneside]{article}

% ---------- Codificación e idioma ----------
% inputenc y fontenc NO se usan con XeLaTeX (maneja UTF-8 de forma nativa)
\usepackage[spanish,es-tabla]{babel}

% ---------- Tipografía estilo plantilla UV ----------
\usepackage{mathptmx}                 % Times-like body text
\usepackage[scaled=0.92]{helvet}      % Helvetica-like (sustituye Gill Sans en títulos)

% ---------- Geometría (acorde a la plantilla UV) ----------
\usepackage[
  letterpaper,
  top=2.5cm,
  bottom=2.5cm,
  left=3cm,
  right=2.5cm,
  headheight=15pt
]{geometry}

% ---------- Interlineado 21 pts ≈ 1.5 ----------
\usepackage{setspace}
\setstretch{1.4}

% ---------- Paquetes auxiliares ----------
\usepackage{graphicx}
\usepackage{xcolor}
\usepackage{tcolorbox}
\tcbuselibrary{skins,breakable}
\usepackage{enumitem}
\usepackage{fancyhdr}
\usepackage{titlesec}
\usepackage{caption}
\usepackage{subcaption}
\usepackage{array}
\usepackage{tabularx}
\usepackage{longtable}
\usepackage{booktabs}
\usepackage{listings}
\usepackage{hyperref}
\usepackage{float}

% ---------- Colores institucionales (azul y verde UV) ----------
\definecolor{uvblue}{RGB}{30,89,150}
\definecolor{uvgreen}{RGB}{80,165,86}
\definecolor{uvgray}{RGB}{90,90,90}
\definecolor{lightbg}{RGB}{245,247,250}

% ---------- Hipervínculos ----------
\hypersetup{
  colorlinks=true,
  linkcolor=uvblue,
  citecolor=uvblue,
  urlcolor=uvblue,
  pdftitle={Práctica 2: Emulador básico de un procesador},
  pdfauthor={Yadid A. Soto Hernández, José de J. Berman Islas, Miguel A. Damián Bumas — Universidad Veracruzana}
}

% ---------- Estilo de títulos (Gill Sans-like + uvblue) ----------
\titleformat{\section}
  {\sffamily\Large\bfseries\color{uvblue}}
  {\thesection.}{0.6em}{}
\titleformat{\subsection}
  {\sffamily\large\bfseries\color{uvblue!85!black}}
  {\thesubsection}{0.5em}{}
\titleformat{\subsubsection}
  {\sffamily\normalsize\bfseries\color{uvgray}}
  {\thesubsubsection}{0.4em}{}
\titlespacing*{\section}{0pt}{1.6em}{0.8em}
\titlespacing*{\subsection}{0pt}{1.2em}{0.5em}
\titlespacing*{\subsubsection}{0pt}{0.9em}{0.4em}

% ---------- Encabezado y pie de página estilo UV ----------
\pagestyle{fancy}
\fancyhf{}
\renewcommand{\headrulewidth}{0pt}
\renewcommand{\footrulewidth}{0.4pt}
\fancyfoot[L]{\footnotesize\color{uvgray}\itshape ``Lis de Veracruz: Arte, Ciencia, Luz''}
\fancyfoot[R]{\footnotesize\color{uvblue}\textbf{www.uv.mx}}
\fancyfoot[C]{\footnotesize\color{uvgray}\thepage}

% ---------- Cajas de información ----------
\newtcolorbox{infobox}[1][]{
  colback=lightbg,
  colframe=uvblue,
  boxrule=0.4pt,
  arc=2pt,
  left=8pt,right=8pt,top=6pt,bottom=6pt,
  fonttitle=\sffamily\bfseries\color{white},
  coltitle=white,
  colbacktitle=uvblue,
  #1
}

\newtcolorbox{requisito}[1][]{
  enhanced,
  colback=lightbg,
  colframe=uvgreen!75!black,
  boxrule=0.4pt,
  arc=2pt,
  left=6pt,right=6pt,top=4pt,bottom=4pt,
  #1
}

% ---------- Listings para código ----------
\lstset{
  basicstyle=\ttfamily\small,
  keywordstyle=\color{uvblue}\bfseries,
  commentstyle=\color{uvgreen!60!black}\itshape,
  stringstyle=\color{red!70!black},
  numbers=left,
  numberstyle=\tiny\color{uvgray},
  numbersep=8pt,
  frame=single,
  framerule=0.3pt,
  rulecolor=\color{uvgray!40},
  backgroundcolor=\color{lightbg},
  breaklines=true,
  breakatwhitespace=true,
  showstringspaces=false,
  tabsize=2,
  captionpos=b,
  xleftmargin=12pt,
  xrightmargin=4pt,
  aboveskip=10pt,
  belowskip=10pt
}

% ---------- Pie de figura ----------
\captionsetup{font=small,labelfont={bf,sf,color=uvblue},labelsep=period,
              format=plain,justification=centering}

% ---------- Comando: figura centrada con marco ----------
\newcommand{\figura}[3][0.85\textwidth]{%
  \begin{figure}[H]
    \centering
    \fbox{\includegraphics[width=#1]{#2}}
    \caption{#3}
  \end{figure}
}

% ---------- Comando: figura SIN marco (capturas oscuras) ----------
\newcommand{\captura}[3][0.8\textwidth]{%
  \begin{figure}[H]
    \centering
    \includegraphics[width=#1]{#2}
    \caption{#3}
  \end{figure}
}

% ---------- Comando: ruta/archivo en monoespaciada con quiebres ----------
\usepackage{seqsplit}
\newcommand{\ruta}[1]{\texttt{\seqsplit{#1}}}

% =============================================================
%                          DOCUMENTO
% =============================================================
\begin{document}

% -------------------------------------------------------------
% PORTADA Y ANTEPORTADA: se agregan aquí desde el editor.
% El equipo las tiene en PDF de la plantilla UV.
% -------------------------------------------------------------

% ============== ÍNDICE ==============
\pagenumbering{roman}
\setcounter{page}{2}
{\sffamily
  \renewcommand{\contentsname}{\color{uvblue}\Large\bfseries Índice}
  \tableofcontents
}
\newpage

\pagenumbering{arabic}
\setcounter{page}{1}

% =============================================================
%                       1. INTRODUCCIÓN
% =============================================================
\section{Introducción}

Los procesadores que dan vida a los dispositivos que usamos a diario se
perciben como una caja negra: se sabe que ejecutan instrucciones, pero rara vez
se tiene la oportunidad de construir uno para observar qué ocurre realmente
dentro. En esta práctica desarrollamos, desde cero, un \textbf{emulador del
procesador MIPS} en Java, con la intención de sustituir la memorización de
conceptos por la experiencia de implementarlos.

\subsection{¿Por qué construir un emulador de procesador desde cero?}

Un procesador comercial real no puede estudiarse pieza por pieza: contiene
millones de transistores, su diseño es propiedad industrial y se construye en
equipos grandes a lo largo de años. Diseñar el nuestro desde cero nos permitió
tomar y ver cada una de las decisiones que en un chip comercial vienen ya
tomadas: el ancho de palabra, el mapa de direcciones, el conjunto de
instrucciones, el manejo de los desbordamientos y la forma en que se actualiza
el contador de programa.

Además, implementar el ciclo de instrucción (búsqueda, decodificación,
ejecución, memoria y escritura) desmitifica el hardware: cada instrucción que
escribe un programador en un lenguaje de alto nivel termina, al final, en
operaciones de la ALU y en movimientos de datos entre registros y memoria. Y
algo importante desde el punto de vista didáctico: en un procesador real un
error cuesta silicio, mientras que en un emulador un error es una prueba
fallida que se puede corregir, repetir y documentar.

\subsection{Qué gana un alumno al saber implementar un emulador de un procesador real}

Quien ha construido un emulador deja de ver la arquitectura de computadoras
como una asignatura de memorización y la entiende como un sistema que se puede
razonar: sabe qué registros intervienen en una suma, qué significa que un dato
desborde y por qué un bucle consume más ciclos que otro. En lo práctico,
adquiere programación de bajo nivel (manejo de bits, registros, memoria y
aritmética con y sin signo), capacidad de leer documentación técnica y
traducirla a código, y disciplina de ingeniería de software: diseño modular,
pruebas unitarias, pruebas de integración y pruebas de sistema.

A esto se suma una ventaja profesional cada vez más valorada: la capacidad de
usar la inteligencia artificial generativa con criterio, es decir, de pedirle
apoyo a la herramienta y de comprobar después que lo que ésta devuelve es
correcto.

\subsection{Razón de implementar un emulador para el procesador MIPS}

Elegimos MIPS de entre las opciones disponibles (RISC-V, x86, M1 y AMD) por
varias razones. Primero, es una arquitectura \textbf{RISC} con sólo \textbf{tres
formatos de instrucción fijos} (R, I y J), de modo que la codificación de 32
bits es regular y predecible, lo que la hace factible de implementar en un curso
sin sacrificar rigor. Segundo, es el procesador de referencia en la literatura
de arquitectura de computadoras (por ejemplo, en los libros de Patterson y
Hennessy), por lo que existe abundante material para resolver dudas. Tercero, su
flujo de control es transparente: el PC se incrementa en una palabra por cada
instrucción y las ramas se calculan con desplazamientos relativos, lo que permite
observar con claridad cómo se modifica el flujo de un programa. Por último, sus
servicios de sistema (\texttt{syscall}) permiten resolver la entrada y salida de
datos sin necesidad de emular periféricos complejos.

% =============================================================
%                  2. DESARROLLO DE LA ACTIVIDAD
% =============================================================
\section{Desarrollo de la actividad}

\subsection{Descripción del emulador}

El emulador se implementó en Java con una arquitectura modular de \textbf{cinco
paquetes grandes}, siguiendo el esqueleto del diagrama de clases que un
compañero del equipo compartió como referencia:

\begin{infobox}[title=Estructura del proyecto]
\begin{lstlisting}[language=bash,caption={Estructura de carpetas del repositorio},captionpos=b]
src/
|-- cpu/          # CPU (ciclo de 5 etapas IF/ID/EX/MEM/WB) y PC
|-- ejecucion/    # EjecutorInstrucciones, UnidadControl, SenalesControl
|-- memoria/      # MemoriaSistema, UnidadAlmacenamiento, ROM, RAM,
|                 # DiscoDuro y Registros
|-- alu/          # ALU y enum OperacionALU (incluye LUI)
|-- instruccion/  # Instruccion, Decodificador, Opcode, Funct,
|                 # FormatoInstruccion, ServicioSyscall, ProgramaBuilder,
|                 # Desensamblador, Binario
|-- ui/           # Emulador (texto) y VentanaEmulador (Swing)
|-- programas/    # Sumatoria, Factorial, Integral, Promedio, Perceptron
`-- pruebas/      # MiniTest, PruebasUnitarias, PruebasIntegracion,
                  # PruebasSistema, EjecutorPruebas
\end{lstlisting}
\end{infobox}

La memoria está dividida en \textbf{segmentos con dirección}: la ROM guarda las
instrucciones en las direcciones 0 a 99 y se bloquea una vez cargado el
programa; la RAM guarda los datos del programa en las direcciones 100 a 199; y
el DiscoDuro conserva información persistente en las direcciones 200 a 299.
Todas las unidades trabajan con una \textbf{palabra de 32 bits} y \textbf{100
celdas}.

El procesador cuenta con \textbf{32 registros de 32 bits}, donde \texttt{\$zero}
siempre lee cero e ignora las escrituras, y con un \textbf{contador de programa
(PC) de 32 bits} que contiene la dirección de la siguiente instrucción a
ejecutar: durante la búsqueda de instrucciones el PC se incrementa para apuntar
a la siguiente, y las ramas (\texttt{beq}, \texttt{bne}, \texttt{j},
\texttt{jal}, \texttt{jr}) escriben en él la dirección destino.

Además, la \textbf{CPU nunca muere}: arranca en la dirección 0, el PC es
\emph{circular} dentro del mapa de memoria, la celda 0 contiene un NOP y existe
un límite de seguridad de 100\,000 ciclos; sólo se detiene con la instrucción
propia \texttt{halt} (opcode \texttt{0x3F}) o con el servicio \texttt{syscall}
10.

\subsubsection{Procesador seleccionado y ALU}

Se seleccionó el procesador \textbf{MIPS} [x] frente a RISC-V, x86, M1 y AMD.
La ALU procesa 32 bits y soporta \textbf{trece operaciones}: ADD, ADDU, SUB,
SUBU, AND, OR, XOR, NOR, SLT, SLTU, SLL, SRL y LUI, además de NOP. ADD y SUB
detectan desbordamiento con signo y detienen la emulación; ADDU y SUBU envuelven
el resultado sin avisar. MIPS no tiene registros de banderas, por lo que las
comparaciones se hacen con \texttt{slt}/\texttt{slti}/\texttt{sltu} seguidas de
\texttt{beq} o \texttt{bne}.

\subsubsection{Ciclo de la CPU e instrucciones soportadas}

Cada instrucción pasa por \textbf{cinco etapas}:

\begin{requisito}
\textbf{IF} (instrucción = ROM[PC]; PC = PC + 1) $\rightarrow$
\textbf{ID} (separar opcode, rs, rt, rd, shamt, funct o inmediato y generar las
señales de control) $\rightarrow$ \textbf{EX} (la ALU calcula) $\rightarrow$
\textbf{MEM} (\texttt{lw} lee la memoria, \texttt{sw} la escribe) $\rightarrow$
\textbf{WB} (el resultado llega al registro destino).
\end{requisito}

El conjunto soportado incluye las instrucciones de formato R (\texttt{add},
\texttt{addu}, \texttt{sub}, \texttt{subu}, \texttt{and}, \texttt{or},
\texttt{xor}, \texttt{nor}, \texttt{slt}, \texttt{sltu}, \texttt{sll},
\texttt{srl}, \texttt{mult}, \texttt{div}, \texttt{mflo}, \texttt{mfhi},
\texttt{jr}, \texttt{syscall}), de formato I (\texttt{addi}, \texttt{addiu},
\texttt{andi}, \texttt{ori}, \texttt{xori}, \texttt{slti}, \texttt{sltiu},
\texttt{lui}, \texttt{lw}, \texttt{sw}, \texttt{beq}, \texttt{bne}) y de formato
J (\texttt{j}, \texttt{jal}), con cinco servicios de \texttt{syscall}: imprimir
entero (1), imprimir cadena (4), leer entero (5), imprimir carácter (11) y
terminar (10). Los programas de usuario se arman con \texttt{ProgramaBuilder},
que emite las instrucciones y resuelve las etiquetas de las ramas, sin pasar
por un archivo de texto ensamblado.

\subsection{Evidencia de ejecución del código}

Las pruebas técnicas se ejecutaron con la orden
\texttt{java -cp out pruebas.EjecutorPruebas} y arrojaron \textbf{117 pruebas
unitarias, 44 pruebas de integración y 40 pruebas de sistema}, todas pasadas y
ninguna fallida. La primera tabla reúne la ejecución de los cinco programas de
usuario y de las suites de pruebas; la segunda muestra la traza de la
ejecución del programa factorial línea por línea, con el resultado que arrojó
el emulador. De cada ejecución se capturó el estado al inicio y al final, que
se muestran en las figuras siguientes.

\begin{longtable}{@{}p{7.2cm} p{7.6cm}@{}}
\caption{Ejecución de los programas de usuario (Paso 4) y de las suites de pruebas}\label{tab:comandos}\\
\toprule
\textbf{Línea de código a ejecutar} & \textbf{Resultado del emulador}\\
\midrule
\endfirsthead
\toprule
\textbf{Línea de código a ejecutar} & \textbf{Resultado del emulador}\\
\midrule
\endhead
\texttt{java -cp out ui.Emulador sumatoria}   & ``Sumatoria de 3 a 10 = 52'' (68 ciclos)\\
\texttt{java -cp out ui.Emulador factorial}   & ``Factorial de 7 = 5040'' (63 ciclos)\\
\texttt{java -cp out ui.Emulador integral}    & ``Integral de x\^{}2 dx desde x=1 hasta x=5 = 30'' (56 ciclos)\\
\texttt{java -cp out ui.Emulador promedio}    & ``Vector: 4 8 15 16 23'' y ``Promedio = 13'' (123 ciclos)\\
\texttt{java -cp out ui.Emulador perceptron}  & ``z = 3\ \ y = 1'' (62 ciclos)\\
\texttt{java -cp out ui.Emulador factorial --traza --estado} & Traza por ciclo con el PC y el código en hexadecimal; volcado de registros y mapa de memoria\\
\texttt{java -cp out pruebas.PruebasUnitarias}    & RESUMEN PRUEBAS UNITARIAS / Pasadas: 117\ \ Fallidas: 0\\
\texttt{java -cp out pruebas.PruebasIntegracion}  & RESUMEN PRUEBAS DE INTEGRACION / Pasadas: 44\ \ Fallidas: 0\\
\texttt{java -cp out pruebas.PruebasSistema}      & RESUMEN PRUEBAS DE SISTEMA / Pasadas: 40\ \ Fallidas: 0\\
\texttt{java -cp out pruebas.EjecutorPruebas}     & 201 pruebas en total: resultado global, todas pasaron\\
\bottomrule
\end{longtable}

\begin{longtable}{@{}p{7.2cm} p{7.6cm}@{}}
\caption{Traza de la ejecución del programa factorial, línea por línea}\label{tab:traza}\\
\toprule
\textbf{Línea de código a ejecutar} & \textbf{Resultado del emulador}\\
\midrule
\endfirsthead
\toprule
\textbf{Línea de código a ejecutar} & \textbf{Resultado del emulador}\\
\midrule
\endhead
\texttt{lw \$t0, 100(\$zero)}\ \ (PC = 0) & \$t0 = Mem[100] = 7\\
\texttt{addiu \$v0, \$zero, 4}\ \ (PC = 1) & \$v0 = 4 (servicio: imprimir cadena)\\
\texttt{addiu \$a0, \$zero, 101}\ \ (PC = 2) & \$a0 = 101 (dirección de la cadena ``Factorial de '')\\
\texttt{syscall}\ \ (PC = 3) & Salida en pantalla: ``Factorial de ''\\
\texttt{addu \$a0, \$t0, \$zero}\ \ (PC = 4) & \$a0 = 7\\
\texttt{addiu \$v0, \$zero, 1}\ \ (PC = 5) & \$v0 = 1 (servicio: imprimir entero)\\
\texttt{syscall}\ \ (PC = 6) & Salida en pantalla: ``7''\\
\texttt{addiu \$v0, \$zero, 4}\ \ (PC = 7) & \$v0 = 4 (servicio: imprimir cadena)\\
\texttt{addiu \$a0, \$zero, 115}\ \ (PC = 8) & \$a0 = 115 (dirección de la cadena `` ='')\\
\texttt{syscall}\ \ (PC = 9) & Salida en pantalla: `` = ''\\
\texttt{addiu \$t1, \$zero, 1}\ \ (PC = 10) & \$t1 = 1 (acumulador del factorial)\\
\texttt{addiu \$t2, \$zero, 1}\ \ (PC = 11) & \$t2 = 1 (contador del bucle)\\
\texttt{slt \$t3, \$t0, \$t2}\ \ (PC = 12, 1.\textsuperscript{a} vez) & \$t3 = 0, porque 7 \textless{} 1 es falso\\
\texttt{bne \$t3, \$zero, salida}\ \ (PC = 13, 1.\textsuperscript{a} vez) & No salta (\$t3 = 0); PC = 14\\
\texttt{mult \$t0, \$t2}\ \ (PC = 14, 1.\textsuperscript{a} vez) & LO:HI = 7 × 1 = 7\\
\texttt{mflo \$t1}\ \ (PC = 15, 1.\textsuperscript{a} vez) & \$t1 = LO = 7\\
\texttt{addiu \$t2, \$t2, 1}\ \ (PC = 16, 1.\textsuperscript{a} vez) & \$t2 = 2\\
\texttt{j 12}\ \ (PC = 17) & PC = 12 (regresa al inicio del bucle)\\
Bucle PC 12 a 17 repetido hasta completar 7 multiplicaciones & \$t1 = 7 × 6 × 5 × 4 × 3 × 2 × 1 = 5040; \$t2 = 8\\
\texttt{slt \$t3, \$t0, \$t2}\ \ (PC = 12, 8.\textsuperscript{a} vez) & \$t3 = 1, porque 7 \textless{} 8 es verdadero\\
\texttt{bne \$t3, \$zero, salida}\ \ (PC = 13, 8.\textsuperscript{a} vez) & Sí salta: PC = 18 (sale del bucle)\\
\texttt{addu \$a0, \$t1, \$zero}\ \ (PC = 18) & \$a0 = 5040\\
\texttt{addiu \$v0, \$zero, 1}\ \ (PC = 19) & \$v0 = 1 (servicio: imprimir entero)\\
\texttt{syscall}\ \ (PC = 20) & Salida en pantalla: ``5040''\\
\texttt{addiu \$v0, \$zero, 4}\ \ (PC = 21) & \$v0 = 4 (servicio: imprimir cadena)\\
\texttt{addiu \$a0, \$zero, 119}\ \ (PC = 22) & \$a0 = 119 (dirección del salto de línea)\\
\texttt{syscall}\ \ (PC = 23) & Salida en pantalla: salto de línea\\
\texttt{addiu \$v0, \$zero, 10}\ \ (PC = 24) & \$v0 = 10 (servicio: terminar el programa)\\
\texttt{halt}\ \ (PC = 25) & La CPU se detiene con HALT; 63 ciclos; salida final: ``Factorial de 7 = 5040''\\
\bottomrule
\end{longtable}

\subsubsection{Capturas de ejecución}

Se muestra, para cada programa, el estado al \textbf{inicio} y al \textbf{fin} de
la ejecución, tal como lo reporta el emulador.

\captura[0.72\textwidth]{imagenes/ev_sumatoria_inicio.png}{Sumatoria: estado inicial (todos los registros en 0 y el mapa de memoria segmentado).}
\captura[0.72\textwidth]{imagenes/ev_sumatoria_fin.png}{Sumatoria: fin de la ejecución; resultado 52 en 68 ciclos.}
\captura[0.72\textwidth]{imagenes/ev_factorial_inicio.png}{Factorial: estado inicial de la emulación.}
\captura[0.72\textwidth]{imagenes/ev_factorial_fin.png}{Factorial: fin de la ejecución; \$t1 = 5040 y LO = 5040 en 63 ciclos.}
\captura[0.72\textwidth]{imagenes/ev_integral_inicio.png}{Integral: estado inicial de la emulación.}
\captura[0.72\textwidth]{imagenes/ev_integral_fin.png}{Integral: fin de la ejecución; resultado 30 en 56 ciclos.}
\captura[0.72\textwidth]{imagenes/ev_promedio_inicio.png}{Promedio: estado inicial de la emulación.}
\captura[0.72\textwidth]{imagenes/ev_promedio_fin.png}{Promedio: fin de la ejecución; promedio 13 en 123 ciclos.}
\captura[0.72\textwidth]{imagenes/ev_perceptron_inicio.png}{Perceptrón: estado inicial de la emulación.}
\captura[0.72\textwidth]{imagenes/ev_perceptron_fin.png}{Perceptrón: fin de la ejecución; z = 3, y = 1 en 62 ciclos.}

\captura[0.72\textwidth]{imagenes/ev_traza_factorial.png}{Factorial con \texttt{--traza}: primeros ciclos con el PC y el código de la instrucción en hexadecimal.}
\captura[0.72\textwidth]{imagenes/ev_pruebas.png}{Ejecución de las 201 pruebas automáticas: 117 unitarias, 44 de integración y 40 de sistema, todas pasadas.}

% =============================================================
%            3. USO DE INTELIGENCIAS ARTIFICIALES
% =============================================================
\section{Uso de inteligencias artificiales para la generación de código}

La inteligencia artificial generativa se utilizó como herramienta de apoyo en
tres frentes: el \textbf{análisis del diseño}, la \textbf{generación de los
planes de prueba} y la \textbf{redacción de los programas en ensamblador MIPS}.
Los prompts fueron escritos por el equipo y ajustados de forma colaborativa; la
respuesta de la herramienta \textbf{nunca se usó tal cual}: se ejecutó contra las
suites de pruebas y se corrigió hasta que el comportamiento fue el esperado.

A continuación se presentan, en orden, los \textbf{trece prompts} utilizados,
quién los formuló y qué parte del proyecto produjo cada uno, junto con la
captura de pantalla de la conversación.

% ---------- Prompt 1 ----------
\subsection{Prompt 1 — Yadid Soto (análisis del sistema de almacenamiento)}

\begin{infobox}[title=Prompt]
\small
«Analiza las siguientes clases de mi proyecto en Java:
\texttt{UnidadAlmacenamiento}, \texttt{RAM}, \texttt{ROM}, \texttt{Registros} y
contador de programa. \texttt{<CÓDIGO>} ¿Qué unidades funcionales faltan para
emular un procesador completo y cómo las reorganizarías en paquetes sin romper
el código existente?»
\end{infobox}

\textbf{Resultado:} arquitectura modular en cinco paquetes: \texttt{cpu},
\texttt{ejecucion}, \texttt{memoria}, \texttt{alu} y
\ruta{instruccion}; la clase \ruta{MemoriaSistema} con el mapa de direcciones
segmentado y el documento \ruta{docs/arquitectura.md}.

\captura[0.62\textwidth]{imagenes/ia/p01_analisis_1.png}{Prompt 1: solicitud de análisis de las clases de almacenamiento y de qué unidades faltan.}
\captura[0.62\textwidth]{imagenes/ia/p01_analisis_2.png}{Prompt 1, captura 2 de 4: unidades funcionales que faltan (núcleo de CPU, interconexión, memoria y E/S).}
\captura[0.62\textwidth]{imagenes/ia/p01_analisis_3.png}{Prompt 1, captura 3 de 4: árbol de paquetes propuesto y cómo migrar sin romper el código.}
\captura[0.62\textwidth]{imagenes/ia/p01_analisis_4.png}{Prompt 1, captura 4 de 4: inyección de dependencias, composición y resumen ejecutivo.}

% ---------- Prompt 2 ----------
\subsection{Prompt 2 — José de Jesús Berman (plan de pruebas unitarias)}

\begin{infobox}[title=Prompt]
\small
«Genera un plan de pruebas para el emulador del procesador MIPS. Considera la
unidad de almacenamiento (rango de direcciones y limpieza), los registros
(\$zero de sólo lectura), la ALU (las operaciones y el desbordamiento), la
decodificación de las instrucciones de formato R, I y J y el comportamiento del
PC.»
\end{infobox}

\textbf{Resultado:} \ruta{docs/pruebas/plan\_pruebas\_unitarias.md} y
\ruta{src/pruebas/PruebasUnitarias.java}: \textbf{117 pruebas} automáticas,
todas pasadas.

\captura[0.62\textwidth]{imagenes/ia/p02_unitarias_1.png}{Prompt 2, captura 1 de 4: rango de direcciones y segmentación.}
\captura[0.62\textwidth]{imagenes/ia/p02_unitarias_2.png}{Prompt 2, captura 2 de 4: limpieza de memoria, \$zero y ALU.}
\captura[0.62\textwidth]{imagenes/ia/p02_unitarias_3.png}{Prompt 2, captura 3 de 4: desbordamiento y decodificación de instrucciones.}
\captura[0.62\textwidth]{imagenes/ia/p02_unitarias_4.png}{Prompt 2, captura 4 de 4: comportamiento del PC y matriz de trazabilidad.}

% ---------- Prompt 3 ----------
\subsection{Prompt 3 — José de Jesús Berman (plan de pruebas de integración)}

\begin{infobox}[title=Prompt]
\small
«Genera un plan de pruebas de integración de código para el emulador de
procesador MIPS. Verifica el ciclo completo (IF/ID/EX/MEM/WB), la conexión
\texttt{MemoriaSistema} $\rightarrow$ \texttt{CPU} $\rightarrow$ \texttt{Registros} y RAM,
el funcionamiento de las ramas (beq, bne, j, jr) y el uso de los syscalls.»
\end{infobox}

\textbf{Resultado:} \ruta{docs/pruebas/plan\_pruebas\_integracion.md} y
\ruta{src/pruebas/PruebasIntegracion.java}: \textbf{44 pruebas}, todas
pasadas.

\captura[0.62\textwidth]{imagenes/ia/p03_integracion_1.png}{Prompt 3, captura 1 de 4: objetivo, alcance y estrategia de pruebas de integración.}
\captura[0.62\textwidth]{imagenes/ia/p03_integracion_2.png}{Prompt 3, captura 2 de 4: entorno y verificación del ciclo completo y conexión memoria-CPU.}
\captura[0.62\textwidth]{imagenes/ia/p03_integracion_3.png}{Prompt 3, captura 3 de 4: funcionamiento de las ramas y uso de los syscalls.}
\captura[0.62\textwidth]{imagenes/ia/p03_integracion_4.png}{Prompt 3, captura 4 de 4: matriz de trazabilidad, herramientas y riesgos.}

% ---------- Prompt 4 ----------
\subsection{Prompt 4 — José de Jesús Berman (plan de pruebas de sistema)}

\begin{infobox}[title=Prompt]
\small
«Genera un plan de pruebas de sistema para el código para el emulador del
procesador MIPS. Ejecuta los cinco programas de usuario (sumatoria, factorial,
integral, promedio y perceptron), comprueba sus salidas exactas, que cada uno
termine con \texttt{halt}/\texttt{syscall} 10 y que el número de ciclos sea
razonable.»
\end{infobox}

\textbf{Resultado:} \ruta{docs/pruebas/plan\_pruebas\_sistema.md} y
\ruta{src/pruebas/PruebasSistema.java}: \textbf{40 pruebas}, todas
pasadas.

\captura[0.62\textwidth]{imagenes/ia/p04_sistema_1.png}{Prompt 4, captura 1 de 4: objetivo, alcance y entorno de pruebas de sistema.}
\captura[0.62\textwidth]{imagenes/ia/p04_sistema_2.png}{Prompt 4, captura 2 de 4: casos TS-01 (sumatoria) a TS-04 (promedio).}
\captura[0.62\textwidth]{imagenes/ia/p04_sistema_3.png}{Prompt 4, captura 3 de 4: caso TS-05 (perceptrón), verificación de halt y ciclos.}
\captura[0.62\textwidth]{imagenes/ia/p04_sistema_4.png}{Prompt 4, captura 4 de 4: criterios globales, registro de resultados y entregables.}

% ---------- Prompt 5 ----------
\subsection{Prompt 5 — Miguel Damián (código de factorial; prompt literal de la práctica)}

\begin{infobox}[title=Prompt]
\small
«A continuación te paso una lista de instrucciones del procesador MIPS que
puedo compilar en mi emulador del procesador MIPS. Genera un código en lenguaje
ensamblador para el procesador MIPS que implemente el factorial de un número.
Incluye la explicación de cada línea de código. Sólo utiliza las instrucciones
que te muestro a continuación. \texttt{<add, addu, sub, subu, and, or, xor, nor,
slt, sltu, sll, srl, mult, div, mflo, mfhi, jr, syscall, addi, addiu, andi, ori,
xori, slti, sltiu, lui, lw, sw, beq, bne, j, jal>}»
\end{infobox}

\textbf{Resultado:} \ruta{src/programas/Factorial.java} (25 instrucciones) y
su explicación en \ruta{docs/programas/factorial\_explicacion.md}; al
ejecutarlo el emulador muestra ``Factorial de 7 = 5040'' (63 ciclos).

\captura[0.62\textwidth]{imagenes/ia/p05_factorial_1.png}{Prompt 5: programa de factorial en ensamblador MIPS y su explicación línea por línea.}

% ---------- Prompt 6 ----------
\subsection{Prompt 6 — Yadid Soto (sumatoria)}

\begin{infobox}[title=Prompt]
\small
«Genera un código en lenguaje ensamblador MIPS que realice la sumatoria de los
números desde el 3 hasta el 10 usando un bucle con beq/bne y slt para el
control. Sólo utiliza las instrucciones de mi emulador: addu, addiu, slt, beq,
bne, j, lw, sw y syscall. Incluye la explicación de cada línea.»
\end{infobox}

\textbf{Resultado:} \ruta{src/programas/Sumatoria.java}; el emulador imprime
``Sumatoria de 3 a 10 = 52'' (68 ciclos).

\captura[0.62\textwidth]{imagenes/ia/p06_sumatoria_1.png}{Prompt 6: sumatoria con bucle y control de flujo, explicada línea por línea.}

% ---------- Prompt 7 ----------
\subsection{Prompt 7 — Miguel Damián (integral definida)}

\begin{infobox}[title=Prompt]
\small
«Genera código ensamblador MIPS que calcule la integral definida de
f(x) = x\textsuperscript{2} desde x = 1 hasta x = 5, aproximada con
rectángulos de ancho 1 (1\textsuperscript{2} + 2\textsuperscript{2} +
3\textsuperscript{2} + 4\textsuperscript{2}). Usa mult y mflo para elevar al
cuadrado, slt/beq para el bucle y syscall para imprimir el resultado con un
mensaje. Sólo instrucciones de mi emulador y explica cada línea.»
\end{infobox}

\textbf{Resultado:} \ruta{src/programas/Integral.java}; el emulador imprime
``Integral de x\^{}2 dx desde x=1 hasta x=5 = 30'' (56 ciclos).

\captura[0.62\textwidth]{imagenes/ia/p07_integral_1.png}{Prompt 7: integral de x\textsuperscript{2} con mult/mflo y bucle controlado con slt/beq.}

% ---------- Prompt 8 ----------
\subsection{Prompt 8 — Yadid Soto (promedio de un vector)}

\begin{infobox}[title=Prompt]
\small
«Genera código ensamblador MIPS que calcule el promedio del vector fijo \{4, 8,
15, 16, 23\} almacenado en RAM: acumula los elementos en un bucle con lw y
addu, después divide entre 5 con div y mflo, e imprime el vector y el promedio
con los syscalls de cadena y entero. Sólo instrucciones soportadas por mi
emulador.»
\end{infobox}

\textbf{Resultado:} \ruta{src/programas/Promedio.java}; el emulador imprime
``Vector: 4 8 15 16 23'' y ``Promedio = 13'' (123 ciclos).

\captura[0.62\textwidth]{imagenes/ia/p08_promedio_1.png}{Prompt 8, captura 1 de 2: acumulación del vector y cálculo del promedio.}
\captura[0.62\textwidth]{imagenes/ia/p08_promedio_2.png}{Prompt 8, captura 2 de 2: impresión del vector y del promedio; notas de instrucciones usadas.}

% ---------- Prompt 9 ----------
\subsection{Prompt 9 — Miguel Damián (perceptrón)}

\begin{infobox}[title=Prompt]
\small
«Genera código ensamblador MIPS que evalúe un perceptrón de 3 entradas con
datos fijos: entradas x = \{1, 1, 0\}, pesos w = \{2, 2, -3\} y bias -1.
Acumula el producto punto con mult/mflo en un bucle, suma el bias, aplica el
umbral con slt contra \$zero y muestra el valor de z y la activación y. Sólo
instrucciones de mi emulador, con la explicación de cada línea.»
\end{infobox}

\textbf{Resultado:} \ruta{src/programas/Perceptron.java}; el emulador imprime
``z = 3\ \ y = 1'' (62 ciclos; z = 2·1 + 2·1 + (−3)·0 − 1 = 3, y = 1 porque
z $\geq$ 0).

\captura[0.62\textwidth]{imagenes/ia/p09_perceptron_1.png}{Prompt 9: perceptrón de 3 entradas con producto punto, bias y umbral.}

% ---------- Prompt 10 ----------
\subsection{Prompt 10 — Yadid Soto (clase ALU)}

\begin{infobox}[title=Prompt]
\small
«Implementa en Java la clase ALU de un procesador MIPS de 32 bits con las
operaciones ADD, ADDU, SUB, SUBU, AND, OR, XOR, NOR, SLT, SLTU, SLL, SRL y LUI,
más NOP. Detecta el desbordamiento con signo en ADD y SUB y detiene el programa
en ese caso, mientras que ADDU y SUBU envuelven el resultado. Devuelve también
el resultado de las comparaciones SLT y SLTU.»
\end{infobox}

\textbf{Resultado:} \ruta{src/alu/ALU.java} y
\ruta{src/alu/OperacionALU.java} (13 operaciones + NOP), probadas con las
pruebas unitarias de la ALU.

\captura[0.62\textwidth]{imagenes/ia/p10_alu_1.png}{Prompt 10, captura 1 de 2: implementación de la ALU de 32 bits en Java.}
\captura[0.62\textwidth]{imagenes/ia/p10_alu_2.png}{Prompt 10, captura 2 de 2: pruebas de la ALU y puntos de diseño clave.}

% ---------- Prompt 11 ----------
\subsection{Prompt 11 — Miguel Damián (ventana gráfica del emulador)}

\begin{infobox}[title=Prompt]
\small
«Genera una ventana en Java (Swing) para mi emulador del procesador MIPS con
estos campos: instrucciones en binario (32 bits por línea, acepta hexadecimal y
comentarios con \#), historial con Ciclo, PC, Instrucción, Operación, señales
de control y PC siguiente, los 32 registros, la memoria por segmentos y el
estado CPU (PC, HI, LO); botones: Ejecutar 1 ciclo, Ejecutar programa, Cargar
en el PC, Cargar ejemplo, Reiniciar CPU y Limpiar historial.»
\end{infobox}

\textbf{Resultado:} \ruta{src/ui/VentanaEmulador.java}, usada para observar el
estado de la CPU ciclo a ciclo y las señales de control.

\captura[0.62\textwidth]{imagenes/ia/p11_ventana_1.png}{Prompt 11, captura 1 de 2: ventana Swing con historial, registros y estado de la CPU.}
\captura[0.62\textwidth]{imagenes/ia/p11_ventana_2.png}{Prompt 11, captura 2 de 2: refresco de registros y memoria, y notas de uso.}

% ---------- Prompt 12 ----------
\subsection{Prompt 12 — José de Jesús Berman (documentación del conjunto de instrucciones)}

\begin{infobox}[title=Prompt]
\small
«Genera la documentación del conjunto de instrucciones del procesador MIPS que
soporta mi emulador: formato R con su campo funct, formato I con su opcode,
formato J, los efectos de cada instrucción sobre los registros en binario de 32
bits y la tabla de servicios de syscall con sus argumentos.»
\end{infobox}

\textbf{Resultado:} \ruta{docs/instrucciones.md}, referencia para escribir y
revisar los cinco programas de usuario.

\captura[0.62\textwidth]{imagenes/ia/p12_documentacion_1.png}{Prompt 12, captura 1 de 4: formatos de instrucción MIPS32 e instrucciones de formato R.}
\captura[0.62\textwidth]{imagenes/ia/p12_documentacion_2.png}{Prompt 12, captura 2 de 4: instrucciones de formato I.}
\captura[0.62\textwidth]{imagenes/ia/p12_documentacion_3.png}{Prompt 12, captura 3 de 4: instrucciones de formato J y tabla de servicios syscall.}
\captura[0.62\textwidth]{imagenes/ia/p12_documentacion_4.png}{Prompt 12, captura 4 de 4: tabla de servicios syscall y resumen de codificación binaria.}

% ---------- Prompt 13 ----------
\subsection{Prompt 13 — Equipo (reestructuración final a cinco paquetes)}

\begin{infobox}[title=Prompt]
\small
«Reorganiza el emulador en 5 paquetes grandes (cpu, ejecucion, memoria, alu,
instruccion) al estilo del diagrama del compañero, con la memoria segmentada
por dirección (ROM texto@0, RAM datos@100, disco@200), una CPU que nunca muere
(arranca en 0, PC circular, memoria en 0 = NOP, límite de 100000 ciclos) y una
instrucción HALT para que la CPU sepa que el programa terminó; añade pruebas y
regenera los diagramas.»
\end{infobox}

\textbf{Resultado:} la estructura actual del proyecto, los diagramas de clases
en \texttt{docs/diagramas/} y las \textbf{201 pruebas} (117 + 44 + 40) que pasan
con \texttt{./probar.sh}.

\captura[0.62\textwidth]{imagenes/ia/p13_reestructuracion_1.png}{Prompt 13, captura 1 de 4: árbol de paquetes propuesto, mapa de memoria y paquete memoria.}
\captura[0.62\textwidth]{imagenes/ia/p13_reestructuracion_2.png}{Prompt 13, captura 2 de 4: paquete instrucción (Opcode, Instrucción y Decodificador).}
\captura[0.62\textwidth]{imagenes/ia/p13_reestructuracion_3.png}{Prompt 13, captura 3 de 4: paquete cpu (CPU con límite de ciclos y HALT).}
\captura[0.62\textwidth]{imagenes/ia/p13_reestructuracion_4.png}{Prompt 13, captura 4 de 4: pruebas JUnit 5 de memoria y cpu.}

\subsection{Reflexión sobre el uso de la inteligencia artificial}

La IA se usó como \textbf{copiloto}: análisis del diseño, planes de pruebas,
programas de usuario, ventana y explicaciones. Las \textbf{decisiones} las tomó
el equipo: elegir MIPS, la segmentación de la memoria por dirección, el
\texttt{HALT} propio (opcode \texttt{0x3F}), la CPU que nunca muere y el límite
de 100\,000 ciclos. La \textbf{validación} también: correr las 201 pruebas y las
capturas de \texttt{salidas/}. Varios resultados iniciales de la IA no compilaban
o no se comportaban como se esperaba (por ejemplo, las señales de control, que
se reconstruyeron con un \emph{Builder} para evitar errores de argumentos
posicionales) y se corrigieron a mano.

En síntesis, la IA ahorra horas, pero sólo es útil si se sabe qué pedir y se
puede comprobar la respuesta con pruebas automatizadas.

% =============================================================
%                  4. REFLEXIONES FINALES
% =============================================================
\section{Reflexiones finales}

Cada integrante del equipo comparte, a continuación, su reflexión personal
sobre la experiencia de haber construido el emulador desde cero.

\subsection{Yadid Asiel Soto Hernández}

Mi parte del proyecto fueron las unidades de almacenamiento, la ALU y los
programas de sumatoria y promedio. Al inicio pensé que programar un procesador
era cosa de expertos, pero me di cuenta de que no se requiere ser un programador
de nivel industrial: basta con un nivel intermedio sólido en Java ---clases,
paquetes, encapsulamiento y algo de herencia, como \texttt{DiscoDuro} que se
deriva de \texttt{UnidadAlmacenamiento}---, manejo de arreglos y comodidad con
números binarios y hexadecimales. Lo difícil no es escribir la primera versión
del código, sino depurar: cuando la salida no coincide con la esperada hay que
rastrear el PC, los registros y la memoria ciclo por ciclo. Ese proceso me
enseñó a razonar con estado, algo que se aplica directo a sistemas embebidos y
a la verificación de hardware. En lo personal, valoro la \textbf{libertad} que
nos dio esta práctica: tuvimos libertad para decidir cómo estructurar cada
módulo y qué camino de diseño tomar, y eso hizo que el proyecto se sintiera
verdaderamente nuestro. También aprendí a usar la inteligencia artificial
con criterio: la IA me ayudó a plantear el análisis inicial y los programas,
pero sólo funcionó cuando supe qué pedir y pude comprobar la respuesta con las
pruebas del emulador.

\subsection{Jose de Jesus Berman Islas}

De mí dependió el sistema de pruebas (unitarias, de integración y de sistema),
la documentación técnica y el \texttt{DiscoDuro}. Esta parte me hizo ver que la
\textbf{verificación} es tan importante como la construcción: escribir las 117
pruebas unitarias, las 44 de integración y las 40 de sistema fue lo que nos
permitió detectar errores que, de otro modo, habrían pasado desapercibidos
hasta la entrega. Aprendí a leer la especificación de MIPS y a traducirla a
casos de prueba concretos, y a reportar \textbf{sin maquillar} exactamente lo
que las pruebas arrojaron, incluidos los fallos que encontramos por el camino.
En el trabajo en equipo, lo que más valoro no fue el reparto de tareas, sino la
integración: acordamos convenciones comunes desde el inicio y corríamos las tres
suites antes de dar por terminada cualquier función. Para mí, los valores que
más quedaron en evidencia fueron la \textbf{honestidad} al presentar los
resultados tal cual, la \textbf{responsabilidad} de entregar mi parte a tiempo
y por el medio establecido, y la \textbf{integridad} de presentar nuestro
propio trabajo, usando la inteligencia artificial únicamente como herramienta
de apoyo y dejando constancia de su uso, como exige esta práctica.

\subsection{Miguel Alejandro Damian Bumas}

Mi aportación fueron los programas de factorial, integral y perceptrón y la
ventana gráfica del emulador. Esta práctica me confirmó que implementar un
procesador conecta todas las capas de la computación: quien entiende el ciclo
de instrucción comprende qué hace realmente un compilador y cómo se comporta la
memoria, y deja de programar a ciegas. Traducir la especificación de MIPS a
ensamblador y verlo correr fue lo que más me gustó, y construir la ventana
(Swing) me permitió observar la CPU ciclo a ciclo, con sus registros y señales
de control. Como recomendación para quienes tomen el curso el próximo año, diría
que vale la pena por dos razones: mejora como programador, porque depurar a
nivel de registros enseña a razonar con estado y con tiempos, y enseña a usar
la IA donde corresponde ---como herramienta de apoyo que ahorra horas, siempre
que se sepa qué pedir y se pueda verificar lo que devuelve---. En lo personal,
me llevo la experiencia de haber trabajado en un equipo que se apoyó
mutuamente, con \textbf{solidaridad} cuando alguno se atascaba y con
\textbf{respeto} por las ideas de los demás, entregando un trabajo del que los
tres nos hacemos responsables.

% =============================================================
%                       5. REFERENCIAS
% =============================================================
\section{Referencias}

\begin{itemize}[leftmargin=1.5em,itemsep=4pt]
  \item Patterson, D. A., \& Hennessy, J. L. (2009). \emph{Computer organization
    and design: The hardware/software interface} (4.\textsuperscript{a} ed.).
    Morgan Kaufmann.
  \item Harris, D. G., \& Harris, S. L. (2012). \emph{Digital design and computer
    architecture: MIPS Edition}. Morgan Kaufmann.
  \item Hamacher, V. C., Vranesic, Z. G., Zaky, S. A., \& Manjikian, N. (2002).
    \emph{Computer organization} (5.\textsuperscript{a} ed.). McGraw-Hill.
  \item MIPS Technologies. (s. f.). \emph{MIPS32 architecture for programmers,
    volume I: Introduction to the MIPS32 architecture}. Recuperado de
    \url{https://www.mips.com/}
  \item Universidad Veracruzana. (2026). \emph{Práctica 2. Emulador básico de un
    procesador} [Instrucciones de la asignatura y plantilla del reporte]. Eminus.
\end{itemize}

\end{document}
